\pagebreak
\subsection{Model Operation}
\label{subsec:monitoring}
Proper handling of AI applications during deployment involves managing different aspects. Technical correctness and desired performance of the AI application is one, but IT security, maintenance and MLOps, workforce acceptance, knowledge management, and legal concerns must also be considered. Lawmakers are very clear in their demands \citep{AiAct, EU2019_Ethic}. Article 61 (1) of the AI Act states \emph{``Providers shall establish and document a post-market monitoring system in a manner
that is proportionate to the nature of the artificial intelligence technologies and the
risks of the high-risk AI system.''} The challenge is to establish specific methods to control assessed risks. Ultimately, decisions need to be made individually for each use case. We give an overview over currently discussed topics and introduce our monitoring card as a de facto checklist. 

\subsubsection{Operating Concept}
In order to use an AI application in productive operation, the operational concept defines the way the application interacts with existing processes. The AI Assessment Catalog specifies this to a detailed description of the used software components and APIs to other environments \citep{poretschkin2021_IAIS_Leitfaden}. As AI falls under the category of software, \cite{winter2021_TueV_Austria} suggest similar IT security precautions, as required for other high-impact software. Additionally, the \cite{ISO27000-Cyber-Security} series concerning Cyber Security and the more recent \cite{ISO27001-ISMS} provide valuable recommendations and best practice examples for establishing an IT security management system (ISMS). Furthermore, it is beneficial to draw upon recommendations from classical software engineering practices when formulating the operating concept. \cite{SCHMIDT2013swengineering} state that \emph{``Concepts of operation descriptions usually address the following:}
\emph{
\begin{itemize}
    \item[1.] Statement of the goals and objectives of the software product.
    \item[2.] Mission statement that expresses the set of services the product provides.
    \item[3.] Strategies, policies, and constraints affecting the product.
    \item[4.] Organizations, activities, and interactions among participants and stakeholders.
    \item[5.] Clear statement of responsibilities associated with product development and sustainment.
    \item[6.] Process identification for distributing, training, and sustainment of the product.
    \item[7.] Milestone decision definitions and authorities.''
    \end{itemize}
    }
\noindent Extending these requirements, the AI Act mandates in Art. 13 the disclosure of the use of AI in outcome determination to users, with the level of detail depending on their function. In all cases, it is essential to explain the reasoning behind AI's decisions and the impact of incoming data features, while also implementing a process to override decisions when necessary. Additionally, for both regular and emergency tasks, a failsafe responsibility must be established. Finally, it is recommended to provide training to employees on the usage of the AI application to enhance acceptance and productivity.

\subsubsection{Monitoring}
To ensure technical correctness, continuous evaluation of the AI application is essential. Depending on the risk level, it is advisable to establish either live monitoring or periodic controls. Live monitoring of the application is recommended to quickly detect anomalies, not only for emergency situations but also to enhance AI performance. The AI Act, as specified in Article 9, mandates the implementation of a `risk management system' that assesses risks during deployment and evaluates ongoing mitigation strategies through an iterative process. The objective is to identify conditions under which the AI system can operate safely. Deviating from this realm of safe operation necessitates the deactivation of the AI application and a transition to a failsafe software system.

\subsubsection{Human-Computer-Interface and Transparency}
For monitoring purposes and insights into the functioning of the AI, a suitable interface is required. AI Act Art. 14 specifies: \emph{``It has to be guaranteed that the potentials, limitations and decisions are transparent without bias.''} It is proposed that the interface provides information at different levels of detail. This way operating staff can be supported in carrying out their tasks with crucial process relevant information by the AI, while data scientists might want to get more insights into training meta data to optimize learning curves. Avoiding bias is a current research topic. \citet{kraus2021erklarbare} point out that an explanation of the algorithm alone is not sufficient for transparency. In accordance with the recommendations set forth by \cite{Unesco2022}  \emph{``Transparency   aims   at   providing   appropriate   information to the respective addressees to enable their understanding and foster trust. Specific to the AI system, transparency  can  enable  people  to  understand  how  each  stage  of  an  AI  system  is  put  in  place,  appropriate  to  the  context  and  sensitivity  of  the  AI  system.  It  may  also  include  insight  into  factors  that  affect  a  specific  prediction  or  decision,  and  whether  or  not  appropriate  assurances  (such  as  safety  or  fairness  measures)  are  in  place. In cases of serious threats of adverse human rights impacts,  transparency  may  also  require  the  sharing  of  code or datasets.''}

\subsubsection{Model Drift}
Model drift summarizes the impact of changes in concept, data, or software and their effect on the model predictions. Concept drift can occur when the underlying correlations of input and output variables change over time, while data drift can occur when the known range is left, and software drift can occur when, for example, the data pipeline is modified and data fields are renamed or units are changed. The effects of changes in the model ecosystem should be closely monitored. That is, to prevent a low model performance based on corrupted or unexpected input data, it is best practice to compare the input data with training data. Metrics such as f-divergence may be used to express differences in the data distribution and calculate limits for how far the live model performance may deviate from the training performance.

\subsubsection{Privacy}
If personal data is collected in the course of the AI deployment, secure storage of this data and deletion periods must be established with the help of a data protection officer. In addition, a process for providing information about personal data must be created in order to process requests from data subjects. The \ac{GDPR} imposes a minimization of personal data collection. Methods such as differential privacy may be used to privatize personal data during collection or procession  \citep{dwork2014algorithmic}.

\subsubsection{MLOps}
While the AI Act imposes no specific requirements for MLOps, except for ensuring the safe operation of AI throughout its lifecycle,  \cite{poretschkin2021_IAIS_Leitfaden} provide several recommendations regarding the implementation of secure MLOps practices.

To ensure uninterrupted operation, a well-regulated MLOps process is of utmost importance. The frequency of regular updates can be determined based on risk and technology impact assessments. When there is a higher risk associated with deviations between the training data and model data, more frequent retraining should be considered. It is crucial to store metadata related to the training process and model, including parameters used, training data, and model performance, to enable reproducibility. This approach relies on the availability of rapidly collected and accurately labeled new data. In situations where such data collection and labeling are not feasible, additional oversight of the AI system by trained personnel, such as data scientists, becomes necessary.

Moreover, based on the identified risks, it is advisable to establish a dedicated testing period for updates to both the model and infrastructure. Smooth transitions necessitate clear delineation of responsibilities for the commit and review processes, as well as deployment. In cases where swift rollback is required, well-documented versioning including data is essential. Additionally, employing a mirrored infrastructure enables updates to be installed seamlessly without causing downtime.

\newpage

\small{
\renewcommand\arraystretch{1.3}
\begin{longtable}{p{0.15\linewidth}|p{0.35\linewidth}p{0.15\linewidth}p{0.35\linewidth}} 
    \toprule
    \multicolumn{4}{c}{\bf{Operation Card}}\\
    \toprule
    \bf{Step}           & \bf{Requirement}   &  \bf{AI Act}   & \bf{References} \\
    \hline
    \hline
    \multirow{4}{=}{Scope and aim of monitoring}
                                    & Describe monitored components & Art 61 (1,3)\\\cline{2-4}
                                    & Assess risks and potential dangers according to Use Case Card & Art 9 (2) \\\cline{2-4}
                                    & List safety measures for risks & Art. 9 (4) &\cite{poretschkin2021_IAIS_Leitfaden} \\
                                    
    \specialrule{.075em}{0em}{0em}
    \multirow{4}{=}{Operating concept}
                                    & Create and document utilisation concept & Art. 13, Art. 17, Art. 16 \\\cline{2-4}
                                    & Plan staff training & Art. 9 (4c) & \cite{BMWK2017industrie} \\\cline{2-4}
                                    & Determine responsibilities & &\cite{poretschkin2021_IAIS_Leitfaden}\\
    \specialrule{.075em}{0em}{0em}
    \multirow{4}{=}{Autonomy of application}
                                    & Document decision-making power of AI & Art. 14, Art. 17 & \cite{poretschkin2021_IAIS_Leitfaden, EU2019_Ethic} \\\cline{2-4}
                                    & Determine process to overrule decisions of the AI & Art. 14 (4) d) e) & \cite{poretschkin2021_IAIS_Leitfaden}\\
    
    \specialrule{.075em}{0em}{0em}
    \multirow{3}{=}{Responsibilities and measures} 
                                    & Document component wise: Assessed risk, control interval,              responsibility, measures for emergency & Art. 9 & \cite{poretschkin2021_IAIS_Leitfaden, EU2019_Ethic, BMWK2017industrie} \\

    \specialrule{.075em}{0em}{0em}
    \multirow{5}{=}{Model performance}
                                    & Monitor input and output & Art. 17 (1) d), Art. 61(2) & \cite{poretschkin2021_IAIS_Leitfaden}\\\cline{2-4}
                                    & Detect drifts in input data & &\cite{poretschkin2021_IAIS_Leitfaden}\\\cline{2-4}
                                    & Document metric in use to monitor model performance & Art. 15 (2) \\
                                    
    \specialrule{.075em}{0em}{0em}
    \multirow{4}{=}{AI interface}
                                    & Establish transparent decision making process & Art. 52 & \cite{poretschkin2021_IAIS_Leitfaden, kraus2021erklarbare} \\\cline{2-4}
                                    & Establish insight in model performance on different levels & Art. 14 & \cite{poretschkin2021_IAIS_Leitfaden} \\
    
    \specialrule{.075em}{0em}{0em}
    \multirow{6}{=}{IT security}
                                    & Document individual access to server rooms & Art. 15 & \cite{winter2021_TueV_Austria, BMWK2019AIandLaw}\\\cline{2-4}
                                    & Set and document needed clearance level for changes to AI/deployment/access regulation & &\cite{winter2021_TueV_Austria, BMWK2019AIandLaw} \\\cline{2-4}
                                    & Establish and audit ISMS & & \cite{ISO27000-Cyber-Security, ISO27001-ISMS}\\

    \specialrule{.075em}{0em}{0em}
    \multirow{4}{=}{Privacy}
                                    & Justify and document use or waiver of a privacy preserving algorithm & Art. 10 (5) & \cite{poretschkin2021_IAIS_Leitfaden, dwork2014algorithmic}\\\cline{2-4}
                                    & If applicable, document privacy algorithm and due changes in the monitoring of output data & &\cite{dwork2014algorithmic} \\ 
    \specialrule{.075em}{0em}{0em}
    \pagebreak
    \specialrule{.075em}{0em}{0em}
    \multirow{6}{=}{MLOps}          
                                    & Establish versioned code repository of AI, training and deployment & &\cite{poretschkin2021_IAIS_Leitfaden} \\\cline{2-4}
                                    & Establish maintenance and update schedule & &\cite{poretschkin2021_IAIS_Leitfaden} \\\cline{2-4}
                                    & Set regulation for the retraining of the AI and decision basis for the replacement of a model & & \cite{poretschkin2021_IAIS_Leitfaden} \\

    \specialrule{.075em}{0em}{0em}
    \multirow{6}{=}{Testing and rollout}
                                    & Determine responsibilities for updates & &\cite{poretschkin2021_IAIS_Leitfaden}\\\cline{2-4}
                                    & Document software tests & Art. 17 (1) d) & \cite{poretschkin2021_IAIS_Leitfaden} \\\cline{2-4}
                                    & Set period of time that an update must function stably before it is transferred to live status \\    
    \specialrule{.075em}{0em}{0em} 
    \multirow{2}{=}{Involvement of
Individuals}
                                    & If feasible, describe involvement of affected individuals & \\     

    \bottomrule
\end{longtable}
}

